%! TEX root = ../note-dqm-transform.tex

\section{Proofs of results in \Cref{sec--intro}}
% \label{sec--prf--intro}

\subsection{Proof of claim in \Cref{eg--non-sigma-finite-pushforward}}
\label{sec--prf--eg--non-sigma-finite-pushforward}

Take any set \(E\) such that \(\nu (E) > 0\).
\(E\) must contain a subset of the form \(A \times \{a\}\) where \(a \in \{0,
1\}\) and \(A\) is a set of positive Lebesgue measure, since otherwise
\(\nu (E) = 0\).
If \(a = 1\), then \(T^{- 1} (A \times \{1\}) = \cup_{y \in A} (- \infty, y]
\times \{y\}\).
Define
\begin{equation*}
  g (x, y) = \ind_{A} (y) \ind_{(- \infty, y]} (x).
\end{equation*}
Then \(g \equiv 1_{T^{- 1} (A \times \{1\})}\).
By the Fubini-Tonelli Theorem,
\begin{equation*}
  \nu (A \times \{1\}) = \mu (T^{- 1} (A \times \{1\})) = \mu \left[ g \right] =
  \int \ind_{A} (y) \cdot \left( \int \ind_{(- \infty, y]} (x) \;
  \mathrm{d} x \right) \; \mathrm{d} y = \infty,
\end{equation*}
since \(\int \ind_{(- \infty, y]} (x) \; \mathrm{d} x = \infty\) for every
\(y \in A\) and \(A\) has positive Lebesgue measure.
By essentially the same argument, the same holds for \(a = 0\), i.e.
\(\nu (A \times \{0\}) = \infty\) by \(\int \ind_{(y, \infty)} (x) \;
\mathrm{d} x = \infty\).
Since \(E \supseteq A \times \{a\}\), \(\nu (E) = \infty\).
Therefore, there cannot be a sequence of measurable sets \(\left\{ E_{n}
\right\}\) such that \(\mathbb{R}^{2} = \cup_{n \in \mathbb{N}} E_{n}\) and
\(\nu \left( E_{n} \right) \in (0, \infty)\) for each \(n
\in \mathbb{N}\).
\qed

\subsection{Proof of \Cref{lem--hellinger-diff-properties}}
\label{sec--prf--lem--hellinger-diff-properties}

\subsubsection{Proof of part \ref{lem--hellinger-diff-properties-q-diff-in-L1}
of \Cref{lem--hellinger-diff-properties}}

From the expression for \(\rho_{\tau, \theta}\) in
\eqref{eqn--diff-def-remainder},
\begin{align*}
  q_{\tau} - q_{\theta} - (\tau - \theta)^{\prime} \dot{q}_{\theta} =
  & \,
  g_{\tau}^{2} - g_{\theta}^{2} - (\tau - \theta)^{\prime} \left( 2
  g_{\theta} \dot{g}_{\theta} \right) \\
  =
  & \, \left( g_{\theta} + (\tau - \theta)^{\prime} \dot{g}_{\theta} +
  \rho_{\tau, \theta} \right)^{2} - g_{\theta}^{2} - 2 g_{\theta} \cdot
  (\tau - \theta)^{\prime} \dot{g}_{\theta} \\
  =
  & \, \left( (\tau - \theta)^{\prime} \dot{g}_{\theta} \right)^{2} +
  \rho_{\tau, \theta}^{2} + 2 g_{\theta} \rho_{\tau, \theta} + 2 (\tau -
  \theta)^{\prime} \dot{g}_{\theta} \rho_{\tau, \theta}.
\end{align*}
By the triangle inequality and the Cauchy-Schwarz inequality,
\begin{equation*}
  \left| q_{\tau} - q_{\theta} - (\tau - \theta)^{\prime} \dot{q}_{\theta}
  \right| \leq |\tau - \theta|^{2} \left| \dot{g}_{\theta} \right|^{2} +
  \rho_{\tau, \theta}^{2} + 2 g_{\theta} \left| \rho_{\tau, \theta}
  \right| + 2 |\tau - \theta| \left| \dot{g}_{\theta} \right| \left|
  \rho_{\tau, \theta} \right|.
\end{equation*}
Hence,
\begin{equation*}
  \begin{split}
  \mu \left[ \left| q_{\tau} - q_{\theta} - (\tau - \theta)^{\prime}
  \dot{q}_{\theta} \right| \right] \leq
  & \, |\tau - \theta|^{2} \mu \left[ \left| \dot{g}_{\theta} \right|^{2}
  \right] + \mu \left[ \rho_{\tau, \theta}^{2} \right] \\
  & + 2 \mu \left[ g_{\theta} \left| \rho_{\tau, \theta} \right|
  \right] + 2 |\tau - \theta| \mu \left[ \left| \dot{g}_{\theta} \right| \left|
  \rho_{\tau, \theta} \right| \right].
  \end{split}
\end{equation*}
By the Cauchy-Schwarz inequality (this time in \(\mathscr{L}_{2} (\mu)\)),
\begin{align*}
  \mu \left[ \left| q_{\tau} - q_{\theta} - (\tau - \theta)^{\prime}
  \dot{q}_{\theta} \right| \right] \leq
  & \, |\tau - \theta|^{2} \mu \left[ \left| \dot{g}_{\theta} \right|^{2}
  \right] + \mu \left[ \rho_{\tau, \theta}^{2} \right] + 2 \sqrt{\mu \left[
  g_{\theta}^{2} \right] \mu \left[ \left| \rho_{\tau, \theta} \right|^{2}
  \right]} \\
  & + 2 |\tau - \theta| \sqrt{\mu \left[ \left| \dot{g}_{\theta} \right|^{2}
  \right] \mu \left[ \left| \rho_{\tau, \theta} \right|^{2} \right]} \\
  =
  & \, O \left( |\tau - \theta|^{2} \right) + o \left( |\tau - \theta|^{2}
  \right) + o \left( |\tau - \theta| \right) + |\tau - \theta| o (|\tau -
  \theta|) \\
  =
  & \, o (|\tau - \theta|).
\end{align*}
\qed

\subsubsection{Proof of part
\ref{lem--hellinger-diff-properties-q-singularity-result} of
\Cref{lem--hellinger-diff-properties}}

Use the remainder \(\rho_{\tau, \theta}\) defined in
\eqref{eqn--diff-def-remainder} to write \(g_{\tau} = g_{\theta} + (\tau -
\theta)^{\prime} \dot{g}_{\theta} + \rho_{\tau, \theta}\).
On the set \(\left\{ g_{\theta} = 0 \right\}\),
\(g_{\tau} = (\tau - \theta)^{\prime} \dot{g}_{\theta} + \rho_{\tau, \theta}\).
Hence,
\begin{equation*}
  \begin{split}
    \mu \left[ g_{\tau}^{2} \mathbf{1} \left\{ g_{\theta} = 0 \right\} \right] =
    & \, \mu \left[ \left( (\tau - \theta)^{\prime} \dot{g}_{\theta} +
    \rho_{\tau, \theta} \right)^{2} \mathbf{1} \left\{ g_{\theta} = 0 \right\}
    \right]
    \leq
    2 |\tau - \theta|^{2} \mu \left[ \left| \dot{g}_{\theta} \right|^{2}
    \right] + 2 \mu \left[ \rho_{\tau, \theta}^{2} \right] \\
    =
    & \, O \left( |\tau - \theta|^{2} \right) + o \left( |\tau - \theta|^{2}
    \right) = O \left( |\tau - \theta|^{2} \right),
  \end{split}
\end{equation*}
which proves \eqref{eqn--hellinger-diff-properties-q-singularity-result-1}.

Now we prove \eqref{eqn--hellinger-diff-properties-q-singularity-result}.
Let \(h = \tau - \theta\).
By assumption \(\theta + h = \tau \in \Theta\) and since \(\theta \in
\mathrm{int} (\Theta)\), for \(|h|\) sufficiently small we have \(\theta - h \in
\Theta\) as well.
Then
\begin{equation*}
  g_{\theta + h} + g_{\theta - h} = \left( h^{\prime} \dot{g}_{\theta} +
  \rho_{\theta + h, \theta} \right) + \left( - h^{\prime} \dot{g}_{\theta} +
  \rho_{\theta - h, \theta} \right) = \rho_{\theta + h, \theta} + \rho_{\theta -
  h, \theta}.
\end{equation*}
By non-negativity, \(0 \leq g_{\tau} \leq g_{\theta + h} + g_{\theta - h}\) and
so by the inequality \((u + v)^{2} \leq 2 \left( u^{2} + v^{2} \right)\),
\(g_{\tau}^{2} \leq \left( g_{\theta + h} + g_{\theta - h}
\right)^{2} \leq 2 \rho_{\theta + h, \theta}^{2} + 2 \rho_{\theta - h,
\theta}^{2}\).
Therefore by \eqref{eqn--diff-def-remainder},
\begin{equation*}
  \mu \left[ g_{\tau}^{2} \cdot \mathbf{1} \left\{ g_{\theta} = 0 \right\}
  \right] \leq 2 \mu \left[ \rho_{\theta + h, \theta}^{2} \right] + 2 \mu \left[
  \rho_{\theta - h, \theta}^{2} \right] = o \left( |h|^{2} \right) = o \left(
  |\tau - \theta|^{2} \right),
\end{equation*}
which establishes \eqref{eqn--hellinger-diff-properties-q-singularity-result}.
\qed

\subsubsection{Proof of part
\ref{lem--hellinger-diff-properties-dotg-singularity-result}
of \Cref{lem--hellinger-diff-properties}}

Use \eqref{eqn--diff-def-remainder} to write
\((\tau - \theta)^{\prime} \dot{g}_{\theta} = g_{\tau} - g_{\theta} -
\rho_{\tau, \theta}\).
On \(\left\{ g_{\theta} = 0 \right\}\), \((\tau - \theta)^{\prime}
\dot{g}_{\theta} = g_{\tau} - \rho_{\tau, \theta}\).
Hence \(\left( (\tau - \theta)^{\prime} \dot{g}_{\theta} \right)^{2} \leq 2
  g_{\tau}^{2} + 2 \rho_{\tau, \theta}^{2}\).
Thus, by \eqref{eqn--diff-def-remainder} and
\eqref{eqn--hellinger-diff-properties-q-singularity-result} in part
\ref{lem--hellinger-diff-properties-q-singularity-result} (now proven),
\begin{equation*}
  \mu \left[ \left( (\tau - \theta)^{\prime} \dot{g}_{\theta} \right)^{2}
  \mathbf{1} \left\{ g_{\theta} = 0 \right\} \right] \leq 2 \mu \left[
  g_{\tau}^{2} \mathbf{1} \left\{ g_{\theta} = 0 \right\} \right] + 2
  \mu \left[ \rho_{\tau, \theta}^{2} \mathbf{1} \left\{ g_{\theta} = 0
  \right\} \right] = o \left( |\tau - \theta|^{2} \right).
\end{equation*}
Then we get that with \(C = \mu \left[ \dot{g}_{\theta}
\dot{g}_{\theta}^{\prime} \mathbf{1} \left\{ g_{\theta} = 0 \right\}
\right]\)
\begin{equation*}
  (\tau - \theta)^{\prime} C (\tau - \theta) \leq \psi \left( |\tau -
  \theta|^{2} \right),
\end{equation*}
for some function \(\psi (\cdot)\) that satisfies \(\psi (w) / |w|^{2} \to 0\)
as \(|w| \to 0\).
We are done if we show that \(C\) is the zero matrix, or equivalently that
\(v^{\prime} C v\) for any \(v \neq \mathbf{0}\).
First note that \(C\) is positive semi-definite and so \(v^{\prime} C v \geq
0\).
Set \(h = \alpha v\) for \(\alpha > 0\).
Then
\begin{equation*}
  \alpha^{2} v^{\prime} C v \leq \psi(\alpha v), \text{ and so }
  v^{\prime} C v \leq \frac{\psi(\alpha v)}{|\alpha v|^{2}} \cdot |v|^{2}.
\end{equation*}
Since \(|v| \neq 0\) and \(\psi(h) = o(|h|^{2})\) as \(|h| \to 0\), taking
\(\alpha \to 0\) sends the right side to zero.
So \(v^{\prime}Cv \leq 0\) for every \(v\).
Combined with \(v^{\prime} C v \geq 0\), we get \(v^{\prime} C v = 0\).
Since \(v\) was arbitrary, \(C\) is the zero matrix, which establishes
\eqref{eqn--hellinger-diff-properties-gdot-singularity-result}.
\qed

\subsection{Proof of \Cref{lem--sigma-finite-meas-abs-bicont-prob-meas}}
\label{sec--prf--lem--sigma-finite-meas-abs-bicont-prob-meas}

Since \(\mu\) is \(\sigma\)-finite, there is a countable
\(\mathcal{B}\)-measurable partition of \(\mathbf{X}\), \(\left\{ E_{n}
\right\}_{n \in \mathcal{N}}\) with \(\mathcal{N} \subseteq \mathbb{N}\), such
that \(\mu \left( E_{n} \right) < \infty\) for each \(n \in \mathcal{N}\).
Since \(\mu\) is positive, we can further restrict the selection to ensure that
\(\mu \left( E_{n} \right) > 0\) for each \(n \in \mathcal{N}\).

\textbf{Case 1:} \(\mathcal{N}\) is finite.
Then \(\mu\) is a finite measure and we can define \(\lambda (A) := \mu (A) /
\mu (\mathbf{X})\) for every \(A \in \mathcal{B}\), which is clearly a
probability measure satisfying \(\lambda \ll \mu\).
By \(\mu (A) = \mu (\mathbf{X}) \cdot \lambda (A)\) and \(\mu (\mathbf{X}) \in
(0, \infty)\), \(\lambda (A) = 0\) implies \(\mu (A) = 0\), and hence \(\mu \ll
\lambda\).

\textbf{Case 2:} \(\mathcal{N}\) is infinite.
Without loss of generality, let \(\mathcal{N} = \mathbb{N}\).
Define
\begin{equation*}
  \lambda (A) := \sum_{n = 1}^{\infty} 2^{- n} \frac{\mu \left( E_{n} \cap A
  \right)}{\mu \left( E_{n} \right)}.
\end{equation*}
Clearly, \(\lambda\) is non-negative, \(\sigma\)-additive and \(\lambda (X) =
1\).
If \(\mu (A) = 0\), then \(\mu \left( E_{n} \cap A \right) = 0\) for every \(n
\in \mathbb{N}\) so that \(\lambda (A) = 0\) and hence, \(\lambda \ll \mu\).
If \(\lambda (A) = 0\), then \(\mu \left( E_{n} \cap A \right) = 0\) for every
\(n \in \mathbb{N}\) so that \(\mu (A) = \sum_{n = 1}^{\infty} \mu \left( E_{n}
\cap A \right) = 0\).
Note that the last claim follows since \(E_{n}\) are pairwise disjoint and
\(\bigcup_{n = 1}^{\infty} E_{n} = X\), which implies that for any \(A \in
\mathcal{B}\), \(E_{n} \cap A\) are pairwise disjoint and \(\bigcup_{n =
1}^{\infty} \left( E_{n} \cap A \right) = A\).
Therefore, \(\mu \ll \lambda\).
\qed

\subsection{Proof of \Cref{lem--hellinger-diff-unaffected-by-dom}}
\label{sec--prf--lem--hellinger-diff-unaffected-by-dom}

\begin{equation*}
  \begin{split}
    \lambda \left[ \left| g_{\ast, \tau} - g_{\ast, \theta} -
    \left( \tau - \theta \right)^{\prime} \dot{g}_{\ast, \theta}
    \right|^{r} \right] =
    & \,
    \lambda \left[ \left| g_{\tau} - g_{\theta} - \left( \tau - \theta
    \right)^{\prime} \dot{g}_{\theta} \right|^{r} f^{r} \right] \\
    =
    & \,
    \mu \left[ \left| g_{\tau} - g_{\theta} - \left( \tau - \theta
    \right)^{\prime} \dot{g}_{\theta} \right|^{r} \right].
  \end{split}
\end{equation*}
The right hand side is \(o \left( \left| \tau - \theta \right|^{r} \right)\) by
differentiability of \(\tau \mapsto g_{\tau}\) at \(\theta\).
\qed

\subsection{Proof of \Cref{lem--density-preservation}}
\label{sec--prf--lem--density-preservation}

Let \(q = \mathrm{d} Q / \mathrm{d} \mu\) and as in
\eqref{eqn--density-preservation}, denote
\begin{equation}
  q_{\ast} = \mu \left[ \frac{\mathrm{d} Q}{\mathrm{d} \mu} \middle| T \right]
  \text{ and } p_{\ast} (t) = \mu \left[ \frac{\mathrm{d}
  Q}{\mathrm{d} \mu} \middle| T = t \right], \text{ so that } q_{\ast} (x) =
  p_{\ast} (T (x)) = \left( p_{\ast} \circ T \right) (x).
  \label{eqn--density-preservation-1}
\end{equation}
\(\nu\) is a probability measure and hence \(\sigma\)-finite.
Both \(P \ll \nu\) and \eqref{eqn--density-preservation} are proven if we show
\begin{equation}
  \nu \left[ p_{\ast} \cdot \ind_{A} \right] = P (A) \quad \text{for
  every } A \in \mathcal{B}^{\prime}.
  \label{eqn--density-preservation-2}
\end{equation}

By the integral change of variables formula%
\footnote{See for example \citet[Theorem 3.6.1]{2007bogachevMeasureTheory}},
\begin{equation*}
  \nu \left[ p_{\ast} \cdot \ind_{A} \right] = (\mu \circ T^{- 1}) \left[
  p_{\ast} \cdot \ind_{A} \right] = \mu \left[ \left( p_{\ast} \cdot
  \ind_{A} \right) \circ T \right] = \mu \left[ \left( p_{\ast} \circ T
  \right) \cdot \left( \ind_{A} \circ T \right) \right].
\end{equation*}
By \(\ind_{A} \circ T \equiv \ind_{T^{- 1} (A)}\) and
\eqref{eqn--density-preservation-1},
\begin{equation*}
  \nu \left[ p_{\ast} \cdot \ind_{A} \right] = \mu \left[ q_{\ast}
  \cdot \ind_{T^{- 1} (A)} \right] = \mu \left[ \mu [q | T] \cdot
  \ind_{T^{- 1} (A)} \right].
\end{equation*}
By the definition of a conditional expectation,
\(\mu \left[ \mu [q | T] \cdot \ind_{T^{- 1} (A)} \right] = \mu \left[ q
\cdot \ind_{T^{- 1} (A)} \right]\) and so,
\begin{equation*}
  \nu \left[ p_{\ast} \cdot \ind_{A} \right] =
  \mu \left[ q \cdot \ind_{T^{- 1} (A)} \right] = Q \left( T^{- 1}
  (A) \right) = Q \circ T^{- 1} (A) =
  P (A),
\end{equation*}
which shows \eqref{eqn--density-preservation-2}.
\qed

\subsection{Proof of \texorpdfstring{\Cref{thm--s-L2-diff}}{Theorem
\ref{thm--s-L2-diff}}}

To show Fr\'echet differentiability, we would like to show that given
\(\dot{s}_{\theta}\) in \eqref{eqn--dots-def},
\begin{equation}
  \nu \left[ r_{\tau, \theta}^{2} \right] = o \left( |\tau - \theta|^{2}
  \right), \text{ where } r \left( t; \tau, \theta \right):= r_{\tau, \theta}
  (t) := s_{\tau} (t) - s_{\theta} (t) - (\tau - \theta)^{\prime}
  \dot{s}_{\theta} (t).
  \label{eqn--s-diff-def}
\end{equation}
Define the following
\begin{equation}
  \dot{s}^{(1)}_{\tau, \theta} (t) = \frac{\mu \left[ 2 g_{\theta}
  \dot{g}_{\theta} \middle| T = t \right]}{s_{\tau} (t) + s_{\theta} (t)}
  \mathbf{1} \left\{ s_{\theta} (t) > 0 \right\}.
  \label{eqn--s-diff-1-def}
\end{equation}
Then,
\begin{equation*}
  r_{\tau, \theta} = s_{\tau} - s_{\theta} - (\tau - \theta)^{\prime}
  \dot{s}_{\tau, \theta}^{(1)} + (\tau - \theta)^{\prime} \dot{s}_{\tau,
  \theta}^{(1)} - (\tau - \theta)^{\prime} \dot{s}_{\theta}.
\end{equation*}
and hence
\begin{equation}
  \begin{split}
    r_{\tau, \theta} =
    & \, r_{\tau, \theta, 1} + r_{\tau, \theta, 2}, \\
    \text{where} \quad r_{\tau, \theta, 1} (t) =
    & \, s_{\tau} (t) - s_{\theta} (t) - (\tau - \theta)^{\prime} \dot{s}_{\tau,
    \theta}^{(1)} (t), \\
    \text{and} \quad r_{\tau, \theta, 2} (t) =
    & \, (\tau - \theta)^{\prime} \dot{s}_{\tau, \theta}^{(1)} (t) - (\tau -
    \theta)^{\prime} \dot{s}_{\theta} (t) \\
    =
    & \, \frac{s_{\theta} (t) - s_{\tau} (t)}{s_{\tau} (t) + s_{\theta} (t)}
    \cdot \frac{\mu \left[ g_{\theta} \cdot (\tau - \theta)^{\prime}
    \dot{g}_{\theta} \middle| T = t \right]}{s_{\theta} (t)} \mathbf{1} \left\{
    s_{\theta} (t) > 0 \right\}.
  \end{split}
  \label{eqn--r-breakdown}
\end{equation}
It suffices to show that
\begin{equation}
  \nu \left[ r_{\tau, \theta, j}^{2} \right] = o \left( |\tau - \theta|^{2}
  \right) \quad \text{for each } j \in \{1, 2\}.
  \label{eqn--rj-oh2}
\end{equation}

\subsubsection{Showing \texorpdfstring{\eqref{eqn--rj-oh2}}{\ref{eqn--rj-oh2}}
for \texorpdfstring{\(j = 1\)}{j = 1}}

By non-negativity of \(s_{\tau}\) and \(s_{\theta}\)
\begin{align*}
  \mathbf{1} \left\{ s_{\theta} (t) > 0 \right\} =
  & \, \mathbf{1} \left\{ s_{\theta} (t) > 0 \right\} \mathbf{1} \left\{
  s_{\tau} (t) + s_{\theta} (t) > 0 \right\} \\
  & +
  \mathbf{1} \left\{ s_{\theta} (t) > 0 \right\} \mathbf{1} \left\{
  s_{\tau} (t) + s_{\theta} (t) = 0 \right\} \\
  =
  & \, \mathbf{1} \left\{ s_{\theta} (t) > 0 \right\} \mathbf{1} \left\{
  s_{\tau} (t) + s_{\theta} (t) > 0 \right\}.
\end{align*}
Thus, rewrite \(\dot{s}^{(1)}_{\tau, \theta}\) in \eqref{eqn--s-diff-1-def} as
\begin{equation*}
  \dot{s}^{(1)}_{\tau, \theta} (t) = \frac{\mu \left[ 2 g_{\theta}
  \dot{g}_{\theta} \middle| T = t \right]}{s_{\tau} (t) +
  s_{\theta} (t)} \mathbf{1} \left\{ s_{\theta} (t) > 0 \right\}
  \mathbf{1} \left\{ s_{\tau} (t) + s_{\theta} (t) > 0 \right\}.
\end{equation*}
By standard properties of conditional expectations,
\begin{equation*}
  \begin{split}
    \dot{s}^{(1)}_{\tau, \theta} (t) =
    & \, \frac{\mu \left[ 2 g_{\theta} \dot{g}_{\theta} \mathbf{1}
    \left\{ s_{\theta} (T) > 0 \right\} \middle| T = t
    \right]}{s_{\tau} (t) + s_{\theta} (t)} \mathbf{1} \left\{
    s_{\tau} (t) + s_{\theta} (t) > 0 \right\} \\
    =
    & \, \frac{\mu \left[ 2 g_{\theta}
    \dot{g}_{\theta} \mathbf{1} \left\{ p_{\theta} (T) > 0 \right\}
    \middle| T = t \right]}{s_{\tau} (t) + s_{\theta} (t)}
    \mathbf{1} \left\{ s_{\tau} (t) + s_{\theta} (t) > 0 \right\},
  \end{split}
\end{equation*}
where the second equality follows
since \(s_{\theta} (T) > 0\) if and only if \(p_{\theta} (T) > 0\).
By \Cref{lem--nonneg-f-equal-cond-exp-pos-prodf},
\(q_{\theta} =
q_{\theta} \mathbf{1} \left\{ p_{\theta} (T) > 0 \right\}\)
\(\mu\)-a.e.
Take square roots to see that
\begin{equation*}
  g_{\theta} \mathbf{1} \left\{ p_{\theta} (T) > 0 \right\} =
  g_{\theta} = g_{\theta} \mathbf{1} \left\{ g_{\theta} > 0
  \right\} \qquad \mu\text{-a.e},
\end{equation*}
where the last equality follows from non-negativity of \(g_{\theta}\).
Thus
\begin{equation*}
  \dot{s}^{(1)}_{\tau, \theta} (t) = \frac{\mu \left[ 2 g_{\theta}
  \mathbf{1} \left\{ g_{\theta} > 0 \right\} \dot{g}_{\theta} \middle| T
  = t \right]}{s_{\tau} (t) + s_{\theta} (t)} \mathbf{1} \left\{
  s_{\tau} (t) + s_{\theta} (t) > 0 \right\}.
\end{equation*}
Let \(B = \left\{ s_{\tau} (t) + s_{\theta} (t) > 0 \right\}\).
Combine with the expression in \eqref{eqn--r-breakdown} to conclude that
\begin{equation*}
  r_{\tau, \theta, 1} = \sqrt{\mu \left[ g_{\tau}^{2} \middle| T = t
  \right]} - \sqrt{\mu \left[ g_{\theta}^{2} \middle| T = t \right]} -
  \frac{\mu \left[ 2 g_{\theta} \mathbf{1} \left\{ g_{\theta} > 0
  \right\} (\tau - \theta)^{\prime} \dot{g}_{\theta} \middle| T = t
  \right]}{\sqrt{\mu \left[ g_{\tau}^{2} \middle| T = t \right]} + \sqrt{\mu
  \left[ g_{\theta}^{2} \middle| T = t \right]}} \mathbf{1}_{B}.
\end{equation*}

By definition, \(g_{\tau}, g_{\theta}, \dot{g}_{\theta} \in
\mathscr{L}_{2} (\mu)\).
Furthermore,
\begin{equation*}
  \begin{split}
    0 \leq \frac{g_{\theta} \mathbf{1} \left\{ g_{\theta} > 0
    \right\}}{g_{\tau} + g_{\theta}}  \leq 1 \quad \text{which
    implies} \quad \frac{2 g_{\theta} \mathbf{1} \left\{ g_{\theta} > 0
    \right\} \dot{g}_{\theta}}{g_{\tau} + g_{\theta}} \in
    \mathscr{L}_{2} (\mu).
  \end{split}
\end{equation*}
Therefore it follows that
\begin{equation*}
  g_{\tau} - g_{\theta} - \frac{2 g_{\theta} \mathbf{1} \left\{ g_{\theta} > 0
  \right\} (\tau - \theta)^{\prime} \dot{g}_{\theta}}{g_{\tau} + g_{\theta}} \in
  \mathscr{L}_{2} (\mu).
\end{equation*}
By \Cref{lem--sqrt-integral-inequality},
\begin{align*}
  r_{\tau, \theta, 1}^{2} (t) \leq
  & \, \mu \left[ \left( g_{\tau} - g_{\theta} - \frac{  2
  g_{\theta} \mathbf{1} \left\{ g_{\theta} > 0 \right\} (\tau - \theta)^{\prime}
  \dot{g}_{\theta}}{g_{\tau} + g_{\theta}} \right)^{2}
  \middle| T = t \right] \\
  =
  & \, \mu \left[ \left( g_{\tau} - g_{\theta} - (\tau - \theta)^{\prime}
  \dot{g}_{\theta} + (\tau - \theta)^{\prime} \dot{g}_{\theta} - \frac{2
  g_{\theta} \mathbf{1} \left\{ g_{\theta} > 0 \right\} (\tau - \theta)^{\prime}
  \dot{g}_{\theta}}{g_{\tau} + g_{\theta}} \right)^{2}
  \middle| T = t \right] \\
  =
  & \, \mu \left[ \left( \rho_{\tau, \theta} + \left( 1 - \frac{2
  g_{\theta} \mathbf{1} \left\{ g_{\theta} > 0 \right\}}{g_{\theta +
  h} + g_{\theta}} \right) (\tau - \theta)^{\prime} \dot{g}_{\theta} \right)^{2}
  \middle| T = t \right] \\
  \leq
  & \, 2 \mu \left[ \left( \rho_{\tau, \theta} \right)^{2} \middle| T = t
  \right] + 2 \mu \left[ \left( 1 - \frac{2 \mathbf{1} \left\{ g_{\theta}
  > 0 \right\} g_{\theta}}{g_{\tau} + g_{\theta}} \right)^{2}
  \left( (\tau - \theta)^{\prime} \dot{g}_{\theta} \right)^{2} \middle| T = t
  \right].
\end{align*}

Split the second sum as follows:
\begin{align*}
  1 - \frac{2 g_{\theta} \mathbf{1} \left\{ g_{\theta} > 0
  \right\}}{g_{\tau} + g_{\theta}} =
  & \, \mathbf{1} \left\{ g_{\theta} > 0 \right\} + \mathbf{1} \left\{
  g_{\theta} = 0 \right\} - \frac{2 g_{\theta} \mathbf{1} \left\{
  g_{\theta} > 0 \right\}}{g_{\tau} + g_{\theta}} \\
  =
  & \, \mathbf{1} \left\{ g_{\theta} > 0 \right\} \left( 1 - \frac{2
  g_{\theta}}{g_{\tau} + g_{\theta}} \right) +
  \mathbf{1} \left\{ g_{\theta} = 0 \right\},
\end{align*}
and so,
\begin{align*}
  \left( 1 - \frac{2 \mathbf{1} \left\{ g_{\theta} > 0
  \right\} g_{\theta}}{g_{\tau} + g_{\theta}} \right)^{2}
  \left( (\tau - \theta)^{\prime} \dot{g}_{\theta} \right)^{2} =
  & \,
  \mathbf{1} \left\{ g_{\theta} > 0 \right\} \left( 1 - \frac{2
  g_{\theta}}{g_{\tau} + g_{\theta}} \right)^{2}
  \left( (\tau - \theta)^{\prime} \dot{g}_{\theta} \right)^{2} \\
  & + \mathbf{1} \left\{ g_{\theta} = 0 \right\} \left( (\tau - \theta)^{\prime}
  \dot{g}_{\theta} \right)^{2} \\
  =
  & \, \mathbf{1} \left\{ g_{\theta} > 0 \right\} \left( \frac{g_{\tau} -
  g_{\theta}}{g_{\tau} + g_{\theta}} \right)^{2} \left( (\tau - \theta)^{\prime}
  \dot{g}_{\theta} \right)^{2} \qquad \mu\text{-a.e},
\end{align*}
by \Cref{lem--hellinger-diff-properties}
\ref{lem--hellinger-diff-properties-dotg-singularity-result}.
Thus,
\begin{align*}
  r_{\tau, \theta, 1}^{2} (t) \leq
  & \, 2 \mu \left[ \left( g_{\tau} - g_{\theta} - (\tau - \theta)^{\prime}
  \dot{g}_{\theta} \right)^{2} \middle| T = t \right] \\
  & + 2 \mu \left[ \mathbf{1} \left\{ g_{\theta} > 0 \right\} \left(
  \frac{g_{\tau} - g_{\theta}}{g_{\tau} + g_{\theta}} \right)^{2} \left( (\tau -
  \theta)^{\prime} \dot{g}_{\theta} \right)^{2} \middle| T = t \right] \\
  \leq
  & \, 2 \mu \left[ \rho_{\tau, \theta}^{2} \middle| T = t \right]
  + 2 |\tau - \theta|^{2} \mu \left[ \mathbf{1} \left\{ g_{\theta} > 0 \right\}
  \left( \frac{g_{\tau} - g_{\theta}}{g_{\tau} + g_{\theta}} \right)^{2} \left|
  \dot{g}_{\theta} \right|^{2} \middle| T = t \right].
\end{align*}
Hence,
\begin{equation*}
  \nu \left[ r_{\tau, \theta, 1}^{2} \right] \leq 2 \left( \mu \left[
  \rho_{\tau, \theta}^{2} \right] + |\tau - \theta|^{2} \mu \left[
  \varphi_{\tau, \theta}^{2} \cdot \left| \dot{g}_{\theta} \right|^{2} \right]
  \right) \quad \text{where} \quad \varphi_{\tau, \theta} = \mathbf{1} \left\{
  g_{\theta} > 0 \right\} \frac{g_{\tau} -
  g_{\theta}}{g_{\tau} + g_{\theta}}.
\end{equation*}
Utilizing this, we note that \(g_{\tau} \overset{\mathscr{L}_{2}
(\mu)}{\to} g_{\theta}\) implies \(g_{\tau} \overset{\mu}{\to}
g_{\theta}\) as \(|\tau - \theta| \to 0\).
By the Continuous Mapping Theorem, \(\varphi_{\tau, \theta} \overset{\mu}{\to}
0\).
Since by definition \(\left| \varphi_{\tau, \theta} \right| \leq 1\), we get
that \(\varphi_{\tau, \theta}^{2} \cdot \left| \dot{g}_{\theta} \right|^{2}
\leq \left| \dot{g}_{\theta} \right|^{2}\).
Since \(\mu\) is a probability measure, by the Dominated Convergence Theorem for
convergence in probability (see for example \citet[Theorem 2.8.5, p.
132]{2007bogachevMeasureTheory}), \(\mu \left[ \varphi_{\tau, \theta}^{2} \cdot
\left| \dot{g}_{\theta} \right|^{2} \right] = o (1)\).
Hence,
\begin{equation*}
  \nu \left[ r_{\tau, \theta, 1}^{2} \right] \leq 2 \left( \mu \left[
  \rho_{\tau, \theta}^{2} \right] + |\tau - \theta|^{2} \mu \left[
  \varphi_{\tau, \theta}^{2} \cdot \left| \dot{g}_{\theta} \right|^{2} \right]
  \right) = 2 \left( o \left( |\tau - \theta|^{2} \right) + |\tau - \theta|^{2}
  o (1) \right) = o \left( |\tau - \theta|^{2} \right).
\end{equation*}

\subsubsection{Showing \texorpdfstring{\eqref{eqn--rj-oh2}}{\ref{eqn--rj-oh2}}
for \texorpdfstring{\(j = 2\)}{j = 2}}

Define
\begin{equation*}
  \psi_{\tau, \theta} (t) = \left| \frac{s_{\tau} (t) -
  s_{\theta} (t)}{s_{\tau} (t) + s_{\theta} (t)} \right|.
  % \label{eqn--psi-def}
\end{equation*}
From the expression in \eqref{eqn--r-breakdown},
\begin{align*}
  \left| r_{\tau, \theta, 2} (t) \right| \leq
  & \, \psi_{\tau, \theta} (t) \left| \frac{\mu \left[ g_{\theta} \cdot (\tau -
  \theta)^{\prime} \dot{g}_{\theta} \middle| T = t \right]}{s_{\theta} (t)}
  \right| \cdot \mathbf{1} \left\{ s_{\theta} (t) > 0 \right\} \\
  \leq
  & \, \psi_{\tau, \theta} (t) \frac{\mu \left[ \left| g_{\theta} \cdot (\tau -
  \theta)^{\prime} \dot{g}_{\theta} \right| \middle| T = t \right]}{s_{\theta}
  (t)} \cdot \mathbf{1} \left\{ s_{\theta} (t) > 0 \right\} \\
  =
  & \, \psi_{\tau, \theta} (t) \frac{\mu \left[ g_{\theta} \cdot
  \left| (\tau - \theta)^{\prime} \dot{g}_{\theta} \right| \middle| T = t
  \right]}{s_{\theta} (t)} \cdot \mathbf{1} \left\{ s_{\theta}
  (t) > 0 \right\} \\
  \leq
  & \, |\tau - \theta| \psi_{\tau, \theta} (t) \frac{\mu \left[ g_{\theta}
  \left| \dot{g}_{\theta} \right| \middle| T = t \right]}{s_{\theta} (t)} \cdot
  \mathbf{1} \left\{ s_{\theta} (t) > 0 \right\}.
\end{align*}
Upon squaring and applying the Cauchy-Schwarz inequality,
{\allowdisplaybreaks
\begin{align*}
  \left| r_{\tau, \theta, 2} (t) \right|^{2}
  =
  & \, |\tau - \theta|^{2} \psi_{\tau, \theta} (t)^{2} \frac{\left( \mu \left[
  g_{\theta} \left| \dot{g}_{\theta} \right| \middle| T = t \right]
  \right)^{2}}{s_{\theta} (t)^{2}} \cdot \mathbf{1} \left\{ s_{\theta} (t) > 0
  \right\} \\
  \leq
  & \, |\tau - \theta|^{2} \psi_{\tau, \theta} (t)^{2} \frac{\mu \left[
  g_{\theta}^{2} \middle| T = t \right] \mu \left[ \left| \dot{g}_{\theta}
  \right|^{2} \middle| T = t \right]}{p_{\theta} (t)} \cdot \mathbf{1} \left\{
  p_{\theta} (t) > 0 \right\} \\
  =
  & \, |\tau - \theta|^{2} \psi_{\tau, \theta} (t)^{2} \frac{p_{\theta} (t) \mu
  \left[ \left| \dot{g}_{\theta} \right|^{2} \middle| T = t \right]}{p_{\theta}
  (t)} \cdot \mathbf{1} \left\{ p_{\theta} (t) > 0 \right\} \\
  =
  & \, |\tau - \theta|^{2} \psi_{\tau, \theta} (t)^{2} \cdot \mathbf{1} \left\{
  p_{\theta} (t) > 0 \right\} \cdot \mu \left[ \left| \dot{g}_{\theta}
  \right|^{2} \middle| T = t \right].
\end{align*}}
Hence,
\begin{equation*}
  \left| r_{\tau, \theta, 2} (t) \right|^{2} \leq |\tau - \theta|^{2}
  \psi_{\tau, \theta} (t)^{2} \cdot \mathbf{1} \left\{ p_{\theta} (t) > 0
  \right\} \zeta_{\theta} (t), \quad \text{where} \quad \zeta_{\theta} (t) = \mu
  \left[ \left| \dot{g}_{\theta} \right|^{2} \middle| T = t \right].
\end{equation*}
By the Law of Iterated Expectations, \(\zeta_{\theta} \in
\mathscr{L}_{1} (\nu)\).
By Fr\'echet differentiability of \(\tau \mapsto p_{\tau}\) at
\(\theta\) in \(\mathscr{L}_{1} (\nu)\), \(p_{\tau} \to
p_{\theta}\) in \(\mathscr{L}_{1} (\nu)\) and therefore also in
\(\nu\)-probability.
Hence, \(\psi_{\tau, \theta} \overset{\nu}{\to} 0\) by the Continuous Mapping
Theorem.
Thus, \(\psi_{\tau, \theta}^{2} \leq 1\), \(\psi_{\tau, \theta}^{2}
\overset{\nu}{\to} 0\) and \(\zeta_{\theta} \in \mathscr{L}_{1} (\nu)\).
By the Dominated Convergence Theorem for convergence in probability (see for
example \citet[Theorem 2.8.5, p. 132]{2007bogachevMeasureTheory})
\begin{equation*}
  \nu \left[ r_{\tau, \theta, 2}^{2} \right] = |\tau - \theta|^{2} \cdot o (1) =
  o \left( |\tau - \theta|^{2} \right).
\end{equation*}

%%% Local Variables:
%%% mode: LaTeX
%%% TeX-master: "../note-dqm-transform"
%%% End:
